%%%%%%%%%%%%%%%%%%%%%%%%%%%%%%%%%%%%%%%%%%%%%%%%%%%%%%%%%%%%%%%%%%%%%%%
%%%% Основные сведения о диссертации                              %%%%
%%%% Заполните этот файл своими данными                           %%%%
%%%%%%%%%%%%%%%%%%%%%%%%%%%%%%%%%%%%%%%%%%%%%%%%%%%%%%%%%%%%%%%%%%%%%%%

%%% Автор диссертации %%%
\newcommand{\thesisAuthorLastName}{Русаков}
\newcommand{\thesisAuthorOtherNames}{Константин Дмитриевич}
\newcommand{\thesisAuthorInitials}{К.\,Д.}
\newcommand{\thesisAuthor}{\thesisAuthorLastName~\thesisAuthorOtherNames}
\newcommand{\thesisAuthorShort}{\thesisAuthorInitials~\thesisAuthorLastName}

%%% Название работы %%%
\newcommand{\thesisTitle}{%
    Комплексная методика реидентификации людей в~информационно-поисковых системах
}

%%% Специальность %%%
\newcommand{\thesisSpecialtyNumber}{2.3.8}
\newcommand{\thesisSpecialtyTitle}{%
    Информатика и~информационные процессы
}

%%% Ученая степень %%%
\newcommand{\thesisDegree}{кандидата технических наук}
\newcommand{\thesisDegreeShort}{канд. техн. наук}

%%% Место и год %%%
\newcommand{\thesisCity}{Москва}
\newcommand{\thesisYear}{2026}

%%% Организация, где выполнена работа %%%
\newcommand{\thesisInOrganization}{%
    Федеральном государственном бюджетном учреждении науки Институте
    проблем управления им.\,В.\,А.\,Трапезникова Российской академии наук
    (ИПУ РАН)
}

%%% Научный руководитель %%%
\newcommand{\supervisorFio}{Гончаренко Владимир Иванович}
\newcommand{\supervisorRegalia}{доктор технических наук, доцент}

%%% Официальные оппоненты %%%
\newcommand{\opponentOneFio}{Фамилия Имя Отчество}
\newcommand{\opponentOneRegalia}{доктор физико-математических наук, профессор}
\newcommand{\opponentOneJobPlace}{Название организации}
\newcommand{\opponentOneJobPost}{должность}

\newcommand{\opponentTwoFio}{Фамилия Имя Отчество}
\newcommand{\opponentTwoRegalia}{кандидат физико-математических наук}
\newcommand{\opponentTwoJobPlace}{Название организации}
\newcommand{\opponentTwoJobPost}{должность}

%%% Ведущая организация (опционально) %%%
\newcommand{\leadingOrganizationTitle}{%
    Федеральное государственное бюджетное образовательное учреждение
    высшего образования <<Название ведущей организации>>
}

%%% Защита %%%
\newcommand{\defenseDate}{<<\_\_>> \_\_\_\_\_\_\_\_\_\_ 2026~г.~в~\_\_ часов}
\newcommand{\defenseCouncilNumber}{24.1.107.03}
\newcommand{\defenseCouncilTitle}{Институте проблем управления им.\,В.\,А.\,Трапезникова Российской академии наук}
\newcommand{\defenseCouncilAddress}{117997, Москва, ул.~Профсоюзная, 65}
\newcommand{\defenseCouncilPhone}{+7~(\_\_\_)~\_\_\_-\_\_-\_\_}

%%% Секретарь диссертационного совета %%%
\newcommand{\defenseSecretaryFio}{Е.\,А.~Барабанова}
\newcommand{\defenseSecretaryRegalia}{д-р техн. наук, доцент}

%%% Библиотека для ознакомления с диссертацией %%%
\newcommand{\synopsisLibrary}{Федерального государственного бюджетного учреждения науки Института проблем управления им.\,В.\,А.\,Трапезникова Российской академии наук и~на~сайте www.ipu.ru}
\newcommand{\synopsisDate}{<<\_\_>> \_\_\_\_\_\_\_\_\_\_ 2026~года}

%%% Ключевые слова для метаданных PDF %%%
\newcommand{\keywords}{%
    ключевое слово 1, ключевое слово 2, ключевое слово 3
}

%%% Общая характеристика работы %%%

% Актуальность темы исследования
\newcommand{\actualityText}{%
        В условиях стремительного роста объёмов визуальных данных, поступающих из распределённых
    сетей камер наблюдения и иных сенсоров, существенно возрастает роль
    информационно--поисковых систем (ИПС), ориентированных на автоматический поиск,
    распознавание, идентификацию и реидентификацию человека в видеопотоках.
    Для практических приложений (безопасность, «умный» город, транспортная инфраструктура,
    корпоративные системы контроля доступа) критичны три взаимосвязанных требования:
    оперативность обработки (режим, близкий к реальному времени), масштабируемость при росте
    числа источников видеоданных и надёжность в условиях шумов и неполноты наблюдений.
    Проблематика активно развивается как в зарубежных работах по person re-identification
    и person search (в частности, L.~Zheng, M.~Ye, R.~Vezzani и др.), так и в отечественных
    исследованиях и научных школах, представленных работами А.\,Ю.~Зайцева,
    В.\,Н.~Сахарова, С.\,А.~Лапшина и др., что подтверждает её высокую научную и прикладную
    значимость.

    Типовая архитектура ИПС визуального поиска включает последовательность процессов:
    детектирование объекта в кадре, формирование признакового представления (эмбеддинга),
    сопоставление с базой и принятие решения о принадлежности наблюдаемого субъекта
    к одной из идентичностей. Практика эксплуатации показывает наличие «каскадного эффекта»:
    ошибка детектирования приводит к деградации качества извлечённых признаков, а далее --- к
    ошибкам поиска и реидентификации. Ситуация усложняется динамичностью сцены
    (изменение ракурса, освещённости, плотности людей, окклюзии) и информационной
    неопределённостью: в отдельные моменты времени часть модальностей может отсутствовать
    (например, лицо недоступно из-за поворота головы, перекрытий или низкого разрешения).

    Современные однофакторные решения, опирающиеся только на лицевые либо только на
    нелицевые признаки, недостаточно устойчивы для ИПС: лицо является наиболее
    дискриминативной модальностью, но часто недоступно, тогда как признаки силуэта
    доступнее, но в изолированном виде менее различительны при большом числе кандидатов.
    Следовательно, актуальной является разработка комплексной методики, ориентированной на
    согласованное улучшение ключевых процессов ИПС --- детектирования, формирования
    устойчивых и дискриминативных эмбеддингов и их адаптивной интеграции при принятии
    решения, --- что позволяет повысить точность и устойчивость реидентификации при
    рациональных вычислительных затратах.
}

% Цель работы
\newcommand{\aimText}{%
    Целью работы является повышение эффективности процессов распознавания
    и~реидентификации людей в~информационно‑поисковых системах путем разработки
    комплексной методики, обеспечивающей повышение точности идентификации
    в режиме, близком к реальному времени при рациональных вычислительных затратах за~счёт
    интеграции лицевых и~силуэтных признаков.
}

% Задачи исследования
\newcommand{\tasksText}{%
    Достижение поставленной цели потребовало решение следующих задач:
    \begin{enumerate}
        \item провести анализ существующих подходов к~детектированию, идентификации
              и~реидентификации личности и~оценить их применимость для
              информационно‑поисковых систем, работающих в~условиях динамической
              сцены и~неполноты данных;
        \item разработать модифицированный критерий обучения нейросетевого детектора
              (YOLOv8) для детектирования связанных объектов <<силуэт--лицо>>,
              включающий дополнительный компонент функции потерь
              $\mathcal{L}_{\text{face}}$ и коэффициент $\lambda_{\text{face}}$,
              а также экспериментально исследовать влияние $\lambda_{\text{face}}$
              на показатели качества детекции лиц при сохранении стабильности
              детекции силуэта;
        \item разработать метод формирования дискриминативных эмбеддингов для
              лица и~силуэта на~основе трансформерной архитектуры и~функции
              потерь ArcFace, обеспечивающий устойчивость представлений
              к~изменениям условий наблюдения;
        \item разработать метод интеграции эмбеддингов лица и~силуэта и~принятия
              решения в~задаче ReID на~основе байесовской модели и~MAP‑критерия;
        \item реализовать программную модель и~провести вычислительные эксперименты
              на~репрезентативных наборах данных (CUHK03, Market‑1501, MARS и~др.),
              выполнить сравнение с~современными подходами и~оценить эффективность
              предложенных решений.
    \end{enumerate}
}

% Научная новизна
\newcommand{\noveltyText}{%
    Научная новизна состоит в~разработке:
    \begin{enumerate}
        \item метода обучения нейросетевого детектора YOLOv8 в задаче детектирования
              связанных объектов <<силуэт--лицо>>, отличающегося модификацией функции потерь
              за счёт введения компонента $\mathcal{L}_{\text{face}}$ и коэффициента
              $\lambda_{\text{face}}$, увеличивающего вклад ошибок детекции лиц,
              локализованных внутри рамки силуэта, что повышает полноту и точность
              формирования ROI лица для последующих этапов идентификации и реидентификации
              (п.~4, п.~13 паспорта специальности 2.3.8);
        \item метода формирования дискриминативных эмбеддингов лица и~силуэта,
              ориентированного на~совместное использование модальностей в~ИПС при
              их неполной доступности и~отличающегося двухканальной схемой с~модально‑специфическими
              ViT‑кодировщиками, обучаемыми по~ArcFace на~общем множестве идентичностей,
              а~также процедурой стандартизации эмбеддингов (включая $L_2$‑нормировку),
              обеспечивающей их пригодность для последующей вероятностной интеграции
              (п.~4, п.~13 паспорта специальности 2.3.8);
        \item метода интеграции эмбеддингов лица и~силуэта и~принятия решения
              в~задаче ReID, основанного на~байесовской модели объединения
              модальностей (эмбеддингов лица и~силуэта) и~многоклассовом
              MAP‑критерии, позволяющего корректно принимать решение при временном
              отсутствии или низком качестве одной из~модальностей; метод отличается
              введением механизма автоматической настройки ключевого гиперпараметра
              комбинирования и~обеспечивает адаптивный вклад модальностей
              в~зависимости от~условий наблюдения (п.~1, п.~13 паспорта
              специальности 2.3.8).
    \end{enumerate}
}

% Практическая значимость
\newcommand{\influenceText}{%
    Практическая значимость состоит в~возможности применения разработанного
    методического и~программного обеспечения в~составе информационно‑поисковых
    систем распознавания и~реидентификации людей в изменяющихся условиях наблюдения (освещение, ракурс, окклюзии, изменение числа людей в~кадре),
    а~также в~системах визуального мониторинга и~безопасности, где требуется
    повышенная устойчивость к~отсутствию/снижению качества отдельных модальностей
    и~рациональное использование вычислительных ресурсов.
}

% Методология и методы исследования
\newcommand{\methodsText}{%
    При проведении исследований использованы методы и~средства компьютерного
    зрения и~машинного обучения, включая глубокие нейросетевые архитектуры
    для детектирования объектов на~изображениях и~в~видеопотоках, методы
    обучения представлений (эмбеддингов) с~применением специализированных
    функций потерь метрического обучения (ArcFace), а~также методы
    вероятностного моделирования и~принятия решений (байесовский подход
    и~MAP‑критерий) для интеграции модальностей.

    Для верификации и~сравнения использованы общепринятые метрики качества
    (Precision/Recall, mAP, ROC/PR‑кривые, Rank‑k и~др.) и~вычислительные
    эксперименты на~открытых наборах данных, а~также сравнение с~современными
    (state‑of‑the‑art) подходами по~метрике mAP.
}

% Положения, выносимые на защиту
\newcommand{\defpositionsText}{%
    Положения, выдвигаемые для защиты, состоят в~том, что разработан(-а):
    \begin{enumerate}
        \item метод обучения нейросетевого детектора YOLOv8 для детектирования связанных
             объектов <<силуэт--лицо>> с модифицированной функцией потерь, включающей
             компонент $\mathcal{L}_{\text{face}}$ и коэффициент $\lambda_{\text{face}}$,
             обеспечивающий улучшение метрик детекции лиц при сохранении стабильности
             детекции класса <<силуэт>>; экспериментально установлено, что значения
             $\lambda_{\text{face}}$ в диапазоне 2,0--2,7 обеспечивают максимальные
             показатели по детекции лиц при отсутствии деградации по классу <<силуэт>>
            (п.~4, п.~13 паспорта специальности 2.3.8);
        \item метод формирования дискриминативных эмбеддингов лица и~силуэта
              на~базе ViT и~ArcFace, обеспечивающий получение устойчивых векторных
              представлений и~повышение качества сравнения и~идентификации объектов
              в~задаче ReID (п.~4, п.~13 паспорта специальности 2.3.8);
        \item метод интеграции эмбеддингов лица и~силуэта на~основе байесовской
              модели и~MAP‑критерия, обеспечивающий повышение качества
              реидентификации по~сравнению с~раздельным использованием модальностей
              и~линейным объединением; по~результатам вычислительных экспериментов
              на~наборе CUHK03 получены значения Precision~95,65\%, Recall~95,54\%,
              F1~94,31\%, Accuracy~93,42\%, а~по~метрике mAP достигнуты значения
              95,65\%~/ 98,3\%~/ 89,2\% для наборов CUHK03~/ Market‑1501~/ MARS
              соответственно (п.~1, п.~13 паспорта специальности 2.3.8).
    \end{enumerate}
}

% Достоверность
\newcommand{\reliabilityText}{%
    Достоверность полученных научных результатов подтверждается согласованностью
    теоретических выводов с~результатами вычислительных экспериментов, корректным
    выбором и~применением общепринятых метрик качества, воспроизводимостью
    экспериментальной процедуры и~сравнением с~современными решениями
    на~стандартных наборах данных (CUHK03, Market‑1501, MARS).
}

% Апробация работы
\newcommand{\probationText}{%
    Основные положения и~результаты диссертационной работы докладывались
    и~обсуждались на~плановых научно‑технических конференциях ИПУ~РАН, а~также
    на~международных и~всероссийских конференциях по~информатике, интеллектуальным
    системам и~компьютерному зрению (в~том числе серии MLSD, ICCT, УБС, МКПУ, ВСПУ).

    Результаты диссертационной работы <<Методика распознавания и~реидентификации
    людей в~информационно‑поисковых системах>> использованы в~деятельности
    ФГБУ~<<НИЦ~<<Институт имени Н.\,Е.\,Жуковского>> в~работах по~автоматическому
    распознаванию и~реидентификации объектов. Применение результатов подтверждено
    Актом о~внедрении (утверждено 10.06.2024). В~акте отмечено, что разработанное
    методическое и~программное обеспечение позволяет эффективно детектировать
    объекты и~преобразовывать их в~численные характеристики (эмбеддинги),
    улучшая точность идентификации, а~также обеспечивая повышение эффективности
    и~сокращение времени на~принятие решения до~5\% при использовании систем
    визуального наблюдения и~безопасности.
}

% Личный вклад
\newcommand{\contributionText}{%
    Основные научные результаты получены лично автором. Постановка задачи
    исследования осуществлялась совместно с~научным руководителем. Автором
    выполнены разработка ключевых алгоритмических решений, реализация программных
    модулей, постановка и~проведение вычислительных экспериментов, анализ
    и~интерпретация результатов.
}

% Публикации
\newcommand{\publicationsText}{%
    Результаты диссертационной работы отражены в~13~научных статьях, в~том числе
    в~2~публикациях в~изданиях из~перечня ВАК при Минобрнауки~РФ, 3~статьях
    из~перечня научных изданий, индексируемых в~международных наукометрических
    базах данных Web~of~Science и/или Scopus, а~также в~7~рецензируемых изданиях.
    Диссертационная работа состоит из~введения, 3~глав, заключения, списка
    литературы и~1~приложения. Основная часть работы изложена на~112~страницах,
    содержит 18~иллюстраций, 4~таблицы. Список цитируемой литературы включает
    109~наименований.
}

% Степень разработанности темы
\newcommand{\degreeText}{%
    Задачи распознавания человека по изображению и видео, включая детектирование,
    идентификацию и реидентификацию (ReID), относятся к числу наиболее активно
    развивающихся направлений цифровой обработки аудиовизуальной информации и распознавания
    образов. Для распознавания лица исторически сформировались фундаментальные подходы
    выделения признаков и снижения размерности (в том числе PCA/«eigenfaces», предложенный
    M.~Turk и A.~Pentland), а также современные методы глубокого обучения и метрического
    обучения. Существенный прогресс в обеспечении дискриминативности эмбеддингов связан с
    развитием функций потерь, формирующих выраженные межклассовые границы в признаковом
    пространстве (например, ArcFace, J.~Deng и соавт.), что обеспечило высокую точность
    сопоставления лицевых представлений в разнообразных условиях наблюдения.

    Реидентификация человека как самостоятельное направление компьютерного зрения активно
    развивается с момента её формализации в мультикамерных сценариях (N.~Gheissari и соавт.,
    2006) и последующего бурного роста методов на основе глубокого обучения. В обзорах и
    исследованиях (A.~Bedagkar--Gala, S.\,K.~Shah; R.~Vezzani и соавт.; D.~Wu и соавт.;
    M.~Ye и соавт.) систематизированы постановки ReID/person search, наборы данных,
    метрики качества (mAP, CMC/Rank, ROC/PR) и основные направления развития:
    двухэтапные и одноэтапные схемы, устойчивые к окклюзиям представления,
    part-based/attention подходы, графовые модели, методы переранжирования и др.
    При этом в реальных ИПС по-прежнему проявляется чувствительность качества к ошибкам
    ранних стадий обработки, а также к изменчивости условий съёмки и внешнего вида объекта.

    Отдельную ветвь составляют многомодальные и многофакторные подходы, объединяющие
    комплементарные источники признаков (лицо, силуэт, атрибуты, контекст), что особенно
    важно при неполной доступности модальностей. Предлагаются схемы объединения
    представлений на уровне векторов и оценок (конкатенация, линейное объединение,
    score-level fusion), а также более сложные механизмы внимания и структурные модели.
    Вероятностные методы и байесовские постановки применяются как в задачах распознавания
    лица (B.~Moghaddam и соавт.; D.~Chen и соавт.), так и в задачах ReID и обучения метрик
    (V.\,E.~Liong и соавт.; W.~Liu и соавт.), обеспечивая интерпретируемость правил решения
    и возможность учитывать неопределённость наблюдений.

    В отечественной научной повестке задача реидентификации и построения систем
    идентификации в сложных условиях наблюдения разрабатывается рядом научных школ,
    включая исследования А.\,Ю.~Зайцева, В.\,Н.~Сахарова, С.\,А.~Лапшина и др., где
    рассматриваются как алгоритмические, так и архитектурные аспекты внедрения методов в
    состав прикладных систем видеонаблюдения и видеоаналитики.

    Несмотря на значительное число работ, сохраняется потребность в методически целостном
    решении, ориентированном на эксплуатационные режимы ИПС: (1) повышающем качество
    детектирования связанных объектов «силуэт--лицо» для снижения каскадного эффекта,
    (2) обеспечивающем формирование компактных и дискриминативных эмбеддингов лица и
    силуэта, сопоставимых в рамках единой процедуры принятия решения, и (3) реализующем
    обоснованную интеграцию модальностей с корректной работой при отсутствии одной из них.
    В диссертационной работе развитие темы реализовано в виде комплексной цепочки
    «детектирование --- формирование эмбеддингов --- вероятностная интеграция» с
    экспериментальной проверкой на репрезентативных наборах данных.
}

% Объект исследования
\newcommand{\objectText}{%
    Объектом исследования являются информационные процессы автоматического
    распознавания и~реидентификации людей в~информационно‑поисковых системах
    в~режиме, близком к~реальному времени.
}

% Предмет исследования
\newcommand{\subjectText}{%
    Предметом исследования являются методы, модели и~программно‑алгоритмические
    средства повышения эффективности распознавания и~реидентификации личности
    на~основе анализа и~интеграции лицевых и~силуэтных признаков.
}

% Теоретическая значимость
\newcommand{\theoreticalText}{%
    Теоретическая значимость заключается в~развитии формального подхода
    к~описанию и~оценке информационных процессов распознавания и~реидентификации
    личности в~условиях информационной неопределённости за~счёт построения
    вероятностной модели интеграции признаковых представлений и~обоснования
    решающего правила на~основе MAP‑критерия.
}

% Соответствие диссертации паспорту специальности
\newcommand{\conformityText}{%
    Диссертационная работа выполнена по~специальности 2.3.8 <<Информатика
    и~информационные процессы>> и~соответствует направлениям исследований
    паспорта специальности:
    \begin{itemize}
        \item п.~1 <<Разработка компьютерных методов и~моделей описания, оценки
              и~оптимизации информационных процессов…>>~--- в~работе разработаны
              компьютерные методы описания и~оценки процесса ReID: математическая
              постановка задачи, вероятностная модель интеграции модальностей
              и~критерий принятия решения (MAP);
        \item п.~4 <<Разработка методов и~технологий цифровой обработки
              аудиовизуальной информации… извлечения требуемой информации
              из… изображений, видео контента>>~--- в~работе разработаны методы
              обработки видеоконтента: детектирование объектов <<силуэт--лицо>>
              и~извлечение эмбеддингов по~ROI;
        \item п.~13 <<Разработка и~применение методов распознавания образов…
              нейро‑сетевых… решающих правил…>>~--- в~работе разработаны
              нейросетевые методы и~решающие правила: модифицированный критерий
              обучения детектора, метод формирования эмбеддингов ViT совместно
              с~ArcFace и~решающее правило MAP для многоклассовой идентификации.
    \end{itemize}
}

